\section{最良近似多項式の理論} % 二章
\label{sec:best_approximation_theory}
  本章では，連続関数の多項式による最良近似問題について，存在，一意性，および特徴について記述する．
  これらの理論的結果は，次章で述べるRemesのアルゴリズムの正当性を保証するものである．
  なお，本章の議論は主に参考文献\cite{seminar-note}に基づく．

  \subsection{近似度と最良近似の存在}
   \label{subsec:approximation_error}
    関数$f(x)$を多項式$P(x)$で近似する．$f(x)$は有界閉区間[a,b]上で連続関数とする．
    また，$n$ 次以下の多項式全体の集合を $\mathcal{P}_n$ とする．

    \begin{definition}[最良近似多項式]
    \label{def:best_approximation_poly}
    区間 $[a, b]$ 上の連続関数 $f \in C[a, b]$ と，$n$ 次以下の多項式 $P \in \mathcal{P}_n$ に対して，近似度（一様ノルム）を次式で定義する．
    \begin{equation}
      \| f - P\|_\infty = \max_{x \in [a, b]} |f(x) - P(x)|.
      \label{eq:uniform_norm}
    \end{equation}
    さらに，近似誤差の下限を
    \begin{equation}
      E_n(f) = \inf_{P \in \mathcal{P}_n} \| f - P \|_\infty
      \label{eq:best_approx_error}
    \end{equation}
    とおく． このとき，$\| f - P^* \|_\infty = E_n(f)$ を満たす多項式 $P^* \in \mathcal{P}_n$ を，$f$ の $n$ 次最良近似多項式と呼ぶ．
    \end{definition}

    このような多項式 $P^*$ が実際に存在することは，$\mathcal{P}_n$ が有限次元空間であることから示すことができる．

    \begin{theorem}[最良近似多項式の存在]
      \label{thm:existence}
      $[a, b]$における任意の連続関数 $f$ に対して，$n$ 次最良近似多項式 $P^* \in \mathcal{P}_n$ が存在する．
      すなわち,最良近似誤差の下限を達成する$P^* \in \mathcal{P}_n$ が存在する
      ．
    \end{theorem}

    \begin{proof}
      下限の定義（式 \ref{eq:best_approx_error}）より，近似誤差が下限値 $E_n(f)$ に収束するような$n$次多項式の列 $\{P^{(\nu)}\}_{\nu=1}^\infty$が存在する． すなわち，
      \[ \lim_{\nu \to \infty} \| f - P_\nu \|_\infty = E_n(f) \]
      である． このとき，三角不等式により
      \begin{equation}
        \| P_\nu \|_\infty \le \| f - P_\nu \|_\infty + \| f \|_\infty
      \end{equation}
      が成り立つ． 十分大きな $\nu$ に対して $\| f - P_\nu \|_\infty \le E_n(f) + 1$ と評価できるため，
      \[ \| P_\nu \|_\infty \le E_n(f) + 1 + \| f \|_\infty \]
      となり，列 $\{P_\nu\}$ は一様有界である．
      $\mathcal{P}_n$ は有限次元ベクトル空間（$\mathbb{R}^{n+1}$ と同型）であるため，ボルツァーノ・ワイエルシュトラスの定理（またはハイネ・ボレルの被覆定理）より，$\{P_\nu\}$ は収束する部分列 $\{P_{\nu_k}\}$ を持つ．
      その極限を $P^* = \lim_{k \to \infty} P_{\nu_k} \in \mathcal{P}_n$ とすれば，ノルムの連続性より
      \begin{equation}
        \| f - P^* \|_\infty = \lim_{k \to \infty} \| f - P_{\nu_k} \|_\infty = E_n(f)
      \end{equation}
      となり，$P^*$ が求める最良近似多項式であることが示される．
    \end{proof}

  \subsection{近似誤差の下限評価}
    \label{subsec:lower_bound_error}
    最良近似多項式の特徴付けを行う準備として，近似誤差の下限評価を与える補題を用意する．
    まず，誤差関数の符号が交互に変化するような点列（交代点列）を定義する．
    \begin{definition}[交代点列]
    \label{def:alternation_points}
      
      区間 $[a, b]$ 内の狭義単調増加列 $a \le \xi_0 < \xi_1 < \dots < \xi_{r-1} \le b$ に対して，誤差 $f(\xi_i) - P(\xi_i)$ が交互に符号を変えるとき，$\xi = (\xi_0, \dots, \xi_{r-1})$ を長さ $r$ の交代点列と呼ぶ．
      長さ $r$ の交代点列全体を $\Xi_r(f, P)$ と表す．すなわち，
      \begin{equation}
        \begin{split}
          \Xi_r(f, P) = \{ \xi = (\xi_0, \dots, \xi_{r-1}) \mid & a \le \xi_0 < \dots < \xi_{r-1} \le b, \ \\
          & (f(\xi_i) - P(\xi_i)) \cdot (f(\xi_{i+1}) - P(\xi_{i+1})) < 0 \}
        \end{split}
      \end{equation}
      である．さらに，この点列上での最小誤差を
      \begin{equation}
        \Phi(f, P, \xi) = \min_{0 \le i \le r-1} |f(\xi_i) - P(\xi_i)|
      \end{equation}
    \end{definition}
    \noindent
    と定義する．
    
    すると，以下の評価が成り立つ．

    \begin{lemma}
      \label{lem:lower_bound}
      $P \in \mathcal{P}_n$ とし，$\xi \in \Xi_{n+2}(f, P)$ ならば，
      \begin{equation}
        \Phi(f, P, \xi) \le E_n(f)
      \end{equation}
      が成り立つ．
    \end{lemma}

    \begin{proof}
      背理法で示す．不等式が成り立たない，すなわち $E_n(f) < \Phi(f, P, \xi)$ と仮定する．
      このとき，定理\ref{thm:existence}より，ある $\tilde{P} \in \mathcal{P}_n$ が存在して
      \begin{equation}
        \| f - \tilde{P} \|_\infty < \Phi(f, P, \xi)
      \end{equation}
      を満たす．ここで，差 $P(x) - \tilde{P}(x)$ を考えると，
      \begin{equation}
        P(\xi_i) - \tilde{P}(\xi_i) = (f(\xi_i) - \tilde{P}(\xi_i)) - (f(\xi_i) - P(\xi_i)) \label{eq:best_approx_error}
      \end{equation}
      である．仮定より
      \begin{equation}
      |f(\xi_i) - \tilde{P}(\xi_i)| \le \| f - \tilde{P} \|_\infty < \Phi(f, P, \xi) \le |f(\xi_i) - P(\xi_i)|
      \end{equation}
      であるため，\eqref{eq:best_approx_error} の符号は右辺の第2項 $-(f(\xi_i) - P(\xi_i))$ の符号に依存している．
      
      $\xi \in \Xi_{n+2}(f, P)$ より $f(\xi_i) - P(\xi_i)$ の符号は交互に入れ替わるため，$P(\xi_i) - \tilde{P}(\xi_i)$ の符号も交互に入れ替わる．
      したがって，連続関数の符号の交代性から（中間値の定理） $P(x) - \tilde{P}(x)$ は各区間 $(\xi_i, \xi_{i+1})$ に少なくとも1つの零点を持つ．
      点は $n+2$ 個あるため区間は $n+1$ 個あり，$P - \tilde{P}$ は少なくとも $n+1$ 個の零点を持つことになる．
      $P - \tilde{P} \in \mathcal{P}_n$ であるため，これが $n+1$ 個の零点を持つならば $P = \tilde{P}$ でなければならない．
      すると $\| f - P \|_\infty < \Phi(f, P, \xi) \le \| f - P \|_\infty$ となり矛盾する．
    \end{proof}

  \label{subsec:alternation_theorem}
  \subsection{チェビシェフの交代定理}

    最良近似多項式の必要十分条件を与えるのが，以下の交代定理である．
    この定理は，Remesアルゴリズムの収束判定条件として本研究の根幹をなすものである．

    \begin{theorem}[交代定理]

      \label{thm:koutaiteiri}
      $f \in C[a, b]$ に対し，$P^* \in \mathcal{P}_n$ が最良近似多項式であるための必要十分条件は，長さ $n+2$ の交代偏差点列 $\xi \in \Xi_{n+2}(f, P^*)$ が存在することである．
      ここで交代偏差点列とは，各点 $\xi_i$ が偏差点（誤差が最大ノルムを達成する点），すなわち
      \begin{equation}
        |f(\xi_i) - P^*(\xi_i)| = \|f - P^*\|_\infty
      \end{equation}
      を満たすような交代点列のことである．
    \end{theorem}

    \begin{proof}
      \textbf{（十分性）}
      長さ $n+2$ の交代偏差点列 $\xi \in \Xi_{n+2}(f, P^*)$ が存在すると仮定する．
      このとき， $\Phi(f, P^*, \xi)$ の定義\ref{def:alternation_points}より
      \[ \Phi(f, P^*, \xi) = \min_{i} |f(\xi_i) - P^*(\xi_i)| = \|f - P^*\|_\infty \]
      である． 補題\ref{lem:lower_bound}より，一般に $\Phi(f, P^*, \xi) \le E_n(f)$ が成り立つため，
      \begin{equation}
        \|f - P^*\|_\infty \le E_n(f)
      \end{equation}
      となる． 一方，最良近似誤差の定義\ref{def:best_approximation_poly}より $E_n(f) \le \|f - P^*\|_\infty$ は常に成り立つため，
      \[ \|f - P^*\|_\infty = E_n(f) \]
      となり，$P^*$ は最良近似多項式であることが示される．

      \textbf{（必要性）}
      背理法により示す． $P^*$ が最良近似多項式であるにもかかわらず，交代偏差点列の長さが $n+2$ 未満であると仮定する．
      すなわち，存在する交代偏差点列の中で点の個数が最大となるものを選び，それを $\xi_0 < \xi_1 < \dots < \xi_m$ とおく．このとき，仮定より $m+1 \le n+1$ である．
      
      このとき，誤差関数 $e(x) = f(x) - P^*(x)$ の符号変化を利用して，最大誤差ノルムを減少させる多項式が存在することを示し，$P^*$ が最良近似多項式であることと矛盾することを導く．

      まず，隣り合う偏差点 $\xi_{i-1}, \xi_i$ の間で，誤差関数の符号が変化する点（零点）を $t_i$ ($i=1, \dots, m$) として選ぶことができる．ここで，補正のための多項式 $K(x)$ を次のように定義する．
      \begin{equation}
        K(x) = \prod_{i=1}^m (t_i - x).
        \label{eq:correction_poly}
      \end{equation}
      仮定より $m \le n$ であるため，$K(x)$ の次数は $n$ 以下であり，$K \in \mathcal{P}_n$ である．
      この $K(x)$ は，各区間 $(t_i, t_{i+1})$ において符号が一定であり，その符号変化のパターンは誤差関数 $e(x)$ と一致する．
      ここで，偏差点 $\xi_i$ における符号を合わせるために，$\sigma = \text{sgn}(e(\xi_0) K(\xi_0))$ とおく．

      今，正の定数 $\lambda > 0$ を用いて新しい多項式 $P_\lambda(x) = P^*(x) + \lambda \sigma K(x)$ を考えると，その誤差は
      \[ f(x) - P_\lambda(x) = (f(x) - P^*(x)) - \lambda \sigma K(x) = e(x) - \lambda \sigma K(x) \]
      となる．
      偏差点 $\xi_i$ を含む近傍では，$e(x)$ と $\sigma K(x)$ が同符号であるため，適切な $\lambda$ を選んで $\lambda \sigma K(x)$ を引くことで，誤差の絶対値 $|e(x)|$ を減少させることができる．      
      一方，偏差点以外の場所では誤差が増大する可能性があるが，$\lambda$ を十分に小さくとることで，増大分を抑え込み，全体として
      \begin{equation}
        \| f - P_\lambda \|_\infty < \| f - P^* \|_\infty
      \end{equation}
      を満たすようにできる．
      これは $P^*$ が最良近似多項式であること（誤差ノルムが最小であること）に矛盾する．
      したがって，最良近似多項式が存在するならば，交代偏差点列の長さは $n+2$ 以上でなければならない．
    \end{proof}
  \label{subsec:uniqueness}
    
  \subsection{一意性}
    \begin{theorem}[最良近似多項式の一意性]
    \label{thm:uniqueness}
      連続関数 $f \in C[a, b]$ に対する $n$ 次最良近似多項式 $P^*$ は一意に定まる．
    \end{theorem}

    \begin{proof}
      背理法により示す． $P^*_1, P^*_2 \in \mathcal{P}_n$ が共に $f$ の最良近似多項式であり，$P^*_1 \neq P^*_2$ であると仮定する．
      このとき，最良近似誤差は $E_n(f) = \|f - P^*_1\|_\infty = \|f - P^*_2\|_\infty$ である．
      ここで，これら2つの多項式の平均 $P_{avg} = \frac{1}{2}(P^*_1 + P^*_2)$ を考えると，三角不等式より
      \begin{equation}
        \|f - P_{avg}\|_\infty = \left\| \frac{1}{2}(f - P^*_1) + \frac{1}{2}(f - P^*_2) \right\|_\infty
        \le \frac{1}{2}\|f - P^*_1\|_\infty + \frac{1}{2}\|f - P^*_2\|_\infty = E_n(f)
      \end{equation}
      が成り立つ． 一方，$E_n(f)$ は誤差の下限であるため $\|f - P_{avg}\|_\infty \ge E_n(f)$ である．
      したがって $\|f - P_{avg}\|_\infty = E_n(f)$ となり，$P_{avg}$ もまた最良近似多項式であることがわかる．

      $P_{avg}$ は最良近似多項式であるため，定理 \ref{thm:koutaiteiri}（交代定理）より，長さ $n+2$ の交代偏差点列 $\xi \in \Xi_{n+2}(f, P_{avg})$ が存在する．
      すなわち，ある符号 $\sigma \in \{1, -1\}$ を用いて，各点 $\xi_i$ ($i=0, \dots, n+1$) において以下が成り立つ．
      \begin{equation}
        f(\xi_i) - P_{avg}(\xi_i) = \sigma (-1)^i E_n(f).
      \end{equation}
      この式の左辺に $P_{avg} = \frac{1}{2}(P^*_1 + P^*_2)$ を代入して変形すると，
      \begin{equation}
        \frac{1}{2}(f(\xi_i) - P^*_1(\xi_i)) + \frac{1}{2}(f(\xi_i) - P^*_2(\xi_i)) = \sigma (-1)^i E_n(f)
        \label{eq:uniq_sum}
      \end{equation}
      となる． ここで，$P^*_1, P^*_2$ は最良近似多項式であるため，任意の点での誤差の絶対値は $E_n(f)$ 以下である．
      すなわち，
      \[ |f(\xi_i) - P^*_1(\xi_i)| \le E_n(f), \quad |f(\xi_i) - P^*_2(\xi_i)| \le E_n(f) \]
      である．
      式 (\ref{eq:uniq_sum}) は，「絶対値が $E_n(f)$ 以下の2つの実数の平均が，ちょうど境界値 $\pm E_n(f)$ になっている」ことを意味する．
      これが成り立つのは，2つの実数が共に境界値と等しい場合のみである．
      したがって，すべての $i=0, \dots, n+1$ において
      \begin{equation}
        f(\xi_i) - P^*_1(\xi_i) = f(\xi_i) - P^*_2(\xi_i) = \sigma (-1)^i E_n(f)
      \end{equation}
      が成立しなければならない． よって
      \[ P^*_1(\xi_i) = P^*_2(\xi_i) \quad (i = 0, 1, \dots, n+1) \]
      となり，2つの $n$ 次多項式が異なる $n+2$ 個の点で一致することになる．
      $n$ 次多項式が $n+1$ 個以上の点で一致すれば恒等的に等しいため，$P^*_1 = P^*_2$ となり，仮定に矛盾する．
      ゆえに最良近似多項式は一意に定まる．
    \end{proof}


\newpage  %% 改ページ ==================================================